% p4_context_engineering_deep.tex
% AI Agent 的核心技术：Context Engineering
% 编译方式: xelatex p4_context_engineering_deep.tex

\documentclass[12pt,a4paper]{article}

% ============================================================
% 字体与中文支持
% ============================================================
\usepackage{fontspec}
\usepackage{xeCJK}
\setCJKmainfont[BoldFont=PingFang SC Semibold]{PingFang SC}
\setCJKsansfont[BoldFont=PingFang SC Semibold]{PingFang SC}
\setCJKmonofont{PingFang SC}
\setmainfont{Palatino}
\setsansfont{Helvetica Neue}
\setmonofont{Menlo}

% ============================================================
% 页面布局
% ============================================================
\usepackage[
  top=2.5cm,
  bottom=2.5cm,
  left=2.5cm,
  right=2.5cm
]{geometry}

% ============================================================
% 常用宏包
% ============================================================
\usepackage{amsmath,amssymb}
\usepackage{graphicx}
\graphicspath{{images/}}
\usepackage{longtable,booktabs,array}
\usepackage{enumitem}
\usepackage{xcolor}
\usepackage{hyperref}
\hypersetup{
  colorlinks=true,
  linkcolor=blue!70!black,
  urlcolor=blue!70!black,
  pdfauthor={整理自李宏毅老师 AI Agent 系列课程},
  pdftitle={AI Agent 的核心技术：Context Engineering},
}

% ============================================================
% 段落与间距
% ============================================================
\usepackage{parskip}
\setlength{\parindent}{0pt}
\setlength{\parskip}{6pt plus 2pt minus 1pt}

% ============================================================
% 引用环境美化
% ============================================================
\usepackage{mdframed}
\newmdenv[
  topline=false,
  bottomline=false,
  rightline=false,
  linewidth=3pt,
  linecolor=blue!40,
  innerleftmargin=12pt,
  innerrightmargin=12pt,
  innertopmargin=8pt,
  innerbottommargin=8pt,
  backgroundcolor=blue!3,
  skipabove=10pt,
  skipbelow=10pt,
]{myquote}

% ============================================================
% 节标题样式
% ============================================================
\usepackage{titlesec}
\titleformat{\section}{\Large\bfseries\sffamily}{}{0em}{}[\vspace{-0.5em}\rule{\textwidth}{0.5pt}]
\titleformat{\subsection}{\large\bfseries\sffamily}{}{0em}{}
\titleformat{\subsubsection}{\normalsize\bfseries\sffamily}{}{0em}{}
\setcounter{secnumdepth}{0}  % 不编号

% ============================================================
% 防止 overfull
% ============================================================
\setlength{\emergencystretch}{3em}
\tolerance=1000

% ============================================================
% 图片插入命令
% ============================================================
\newcommand{\keyslide}[2]{%
  \begin{figure}[htbp]
    \centering
    \includegraphics[width=0.92\textwidth]{#1}
    \caption{#2}
  \end{figure}
}

% ============================================================
\begin{document}

% ============================================================
% 标题页
% ============================================================
\begin{titlepage}
  \centering
  \vspace*{3cm}
  {\Huge\bfseries\sffamily AI Agent 的核心技术：Context Engineering\par}
  \vspace{1.5cm}
  {\Large 整理自李宏毅老师 AI Agent 系列课程第 4 集\par}
  \vspace{2cm}
  {\large\sffamily 挑选 \textbullet{} 压缩 \textbullet{} Multi-Agent\par}
  \vfill
  {\large \today\par}
\end{titlepage}

\newpage
\tableofcontents
\newpage

% ============================================================
\section{课程概览}
% ============================================================

本集课程探讨在 AI Agent 时代，让 Agent 成功运行的关键技术——\textbf{Context Engineering}。课程大纲如下：

\begin{enumerate}[leftmargin=2em]
  \item \textbf{Context Engineering 的概念与演进}
  \item \textbf{一个完整的 Context 应该包含什么？}
  \item \textbf{运行 AI Agent 的挑战：过长的上下文}
  \item \textbf{Context Engineering 的三大套路：选择、压缩、Multi-Agent}
\end{enumerate}

% ============================================================
\section{一、Context Engineering 的概念与演进}
% ============================================================

\subsection{什么是 Context Engineering？}

语言模型本质上是在做"文字接龙"（根据输入 $X$ 预测输出 $f(X)$）。如果结果不如预期，我们有两个方向可以调整：
\begin{itemize}[leftmargin=2em]
  \item \textbf{改变模型本身（参数）}：这就是"训练"（Training / Learning）。
  \item \textbf{改变输入（Prompt/Context）}：对于多数闭源大模型，我们无法改变其参数，只能通过给出更合适的输入来引导其输出。
\end{itemize}

在这堂课中，\textbf{没有任何模型被训练，我们只"训练"人类}，学习如何提供更好的上下文。

\subsection{Context Engineering vs. Prompt Engineering}

这两个词本质上指涉相同的概念——把语言模型的输入弄好。但它们关注的重点随时代发生了变化：

\begin{itemize}[leftmargin=2em]
  \item \textbf{Prompt Engineering 时代}：
  在早期（如 GPT-3 时代），语言模型能力较弱，人们热衷于寻找\textbf{"神奇咒语"}来提升模型表现。
  \begin{itemize}
    \item 示例 1：在 prompt 后面加上一连串毫无意义的 "Waves waves waves"，竟然能让 GPT-3 的输出变长（出自论文：\href{https://arxiv.org/abs/2206.03931}{Learning to Generate Prompts for Dialogue Generation through Reinforcement Learning}）。
    \item 示例 2：加入 "Let's think step by step"（一步步思考）来提升数学推理能力（出自论文：\href{https://arxiv.org/abs/2205.11916}{Large Language Models are Zero-Shot Reasoners}）。
    \item 其他咒语：深呼吸、给你小费、关系到世界和平等。
  \end{itemize}
  但随着模型能力增强，它们本就该"使尽全力"做到最好，神奇咒语的加成效果越来越低。

  \item \textbf{Context Engineering 时代}：
  现在人们关注的是：如何\textbf{自动化地管理语言模型的输入}。特别是当 AI Agent 需要多步运行、自主调用工具时，如何保持其上下文的高效与整洁。
\end{itemize}

% ============================================================
\section{二、一个完整的 Context 应该包含什么？}
% ============================================================

现代语言模型的输入（Context）不再只是一两句简单的问句，它通常极其庞大，包含以下几类信息：

\subsection{1. User Prompt (使用者的指令)}

不仅包含任务目标，还要包含\textbf{详细的限制条件和背景情境}。
\begin{itemize}[leftmargin=2em]
  \item \textbf{背景情境的影响}：同样问"水中是什么动物"，如果补充"这是在曼谷运河拍的"，模型就能准确判断那是"水巨蜥"而不是"鳄鱼"。
  \item \textbf{范例 (In-Context Learning)}：提供少量的输入输出范例，能极大提升模型对特定格式或陌生任务的理解力。例如 Gemini 1.5 的技术报告中，仅通过在 Context 中放入一本极其冷门的卡拉芒语（Kalamang）教科书，模型就能获得翻译卡拉芒语的能力。这完全是靠上下文实现的"学习"，模型参数并未改变。
\end{itemize}

\subsection{2. System Prompt (系统提示词)}

开发平台为模型预设的行为准则。

\keyslide{../key_slides/p4_slide_0040.jpg}{System Prompt}

以 Claude Opus 4.1 为例，它的 System Prompt 长达 2500 字，其中规定了：
\begin{itemize}[leftmargin=2em]
  \item 它是谁、今天是几月几号（解答了"模型为什么知道日期"的疑惑）。
  \item 不要教人制造核武或合成毒药。
  \item 回答时不要老是用 "Good question" 开头。
  \item 如果被人类纠正，不要马上承认错误，要仔细思考。
\end{itemize}

\subsection{3. 记忆与外部资讯}

\begin{itemize}[leftmargin=2em]
  \item \textbf{对话历史记录（短期记忆）}：同一对话窗口中的内容。
  \item \textbf{长期记忆}：如 ChatGPT 的 Memory 功能，它会在背景中根据过去的对话自动生成摘要记录。
  \item \textbf{外部搜索资讯 (RAG)}：通过搜索引擎获取最新信息，加入上下文中以弥补模型知识的滞后。
\end{itemize}

\subsection{4. 工具使用 (Tool Use) 的说明与输出}

如果要让模型使用工具，必须把"工具说明书"塞进 Context 中。

\keyslide{../key_slides/p4_slide_0060.jpg}{Tool Use}

\textbf{模型是如何使用工具的？}
\begin{enumerate}[leftmargin=2em]
  \item 告知模型可用工具的格式（如放在 \texttt{<tool>} 和 \texttt{</tool>} 之间）。
  \item 模型进行文字接龙，生成了一段调用工具的文本。
  \item \textbf{关键点}：文字本身不会执行工具。需要一个外部的小程序（如 Python 中的 \texttt{eval()} 函式），将这段文字提取出来执行。
  \item 将工具返回的结果放到 Context 中。
  \item 模型再根据工具的结果继续接龙出最终答案给使用者。
\end{enumerate}

\keyslide{../key_slides/p4_slide_0100.jpg}{Colab Code}

\textbf{更强大的工具：Computer Use}

像 Claude 或 ChatGPT 代理模式，可以通过截取屏幕画面作为观察（Observation），然后输出移动鼠标或敲击键盘的坐标和指令，实现像人类一样订高铁票或浏览网页。

\subsection{5. 模型自己的思考过程 (Reasoning)}

如 OpenAI 的 o1 或 DeepSeek-R1，它们在给出答案前，会在脑内演练"小剧场"（提出假设、自我验证、修改路线）。这些思考过程也是其 Context 的一部分。

% ============================================================
\section{三、运行 AI Agent 的挑战：输入过长}
% ============================================================

AI Agent 是一个在\textbf{环境}中不断循环的智能体：
\begin{myquote}
观察 (Observation) $\rightarrow$ 行动 (Action) $\rightarrow$ 环境改变 $\rightarrow$ 新的观察 $\rightarrow$ \dots
\end{myquote}

由于 Agent 需要长时间运行，其互动的历史记录会越来越长。

\keyslide{../key_slides/p4_slide_0160.jpg}{Evolution of Context Windows}

虽然现在模型的 Context Window 越来越长（GPT-4 支持 3 万 token，Claude 支持 10 万，Gemini 支持百万，Llama-4 甚至达到千万 token），但这并不意味着"把所有东西都塞给模型"就万事大吉了。

\textbf{过长的上下文会导致灾难：}

\begin{enumerate}[leftmargin=2em]
  \item \textbf{Lost in the Middle (中间迷失)}：研究表明（\href{https://arxiv.org/abs/2307.03172}{Lost in the Middle}），模型通常只记得上下文的开头和结尾。如果关键答案藏在长文档的中间，模型的正确率会断崖式下降。
  \item \textbf{多轮对话中的迷失}：有论文（\href{https://arxiv.org/abs/2505.06120}{LLMs Get Lost In Multi-Turn Conversation}）发现，将一个复杂问题拆解为多个子步骤进行"挤牙膏式"的交互，模型的表现反而比一次性把需求说清楚更差、更不稳定。
  \item \textbf{Context Rot (上下文腐烂)}：
  \keyslide{../key_slides/p4_slide_0170.jpg}{Context Rot}
  即使是最简单的"复读机"任务，当输入长度达到数万 token 时，那些号称支持百万级窗口的模型也会开始胡言乱语（\href{https://arxiv.org/abs/2512.24601}{Recursive Language Models}）。
\end{enumerate}

\textbf{结论}：Context Engineering 的核心目标就是一句话——\textbf{避免塞爆 Context，只放需要的，清掉不需要的。}

% ============================================================
\section{四、Context Engineering 的三大基本套路}
% ============================================================

\subsection{1. 挑选 (Selection)}

不要把所有的资讯都丢给模型，而是利用 RAG 等技术精准挑选。

\keyslide{../key_slides/p4_slide_0180.jpg}{挑选机制}

\begin{itemize}[leftmargin=2em]
  \item \textbf{挑选文章与句子}：先用搜索引擎搜文章，再用一个跑得极快的小模型（如 \href{https://arxiv.org/abs/2501.16214}{Provence 论文} 提出的方案）挑选出真正相关的句子，最后才喂给大模型。
  \keyslide{../key_slides/p4_slide_0178.jpg}{挑选需要的内容 - Provence}
  \item \textbf{挑选工具 (Tool RAG)}：不要把几千个工具的说明书全塞给模型，根据用户的当前需求，只检索相关的 3~5 个工具放入 Context 中。
  \item \textbf{挑选记忆 (Memory RAG)}：在早期的\href{https://arxiv.org/abs/2304.03442}{斯坦福虚拟小镇 (Generative Agents)}中，Agent 的记忆是极其琐碎的。系统通过给记忆打分（根据近期程度、重要性、相关性），只提取最高分的记忆给大模型。
\end{itemize}

\keyslide{../key_slides/p4_slide_0185.jpg}{StreamBench}
\keyslide{../key_slides/p4_slide_0190.jpg}{StreamBench 表现}

\textbf{挑选历史经验的诡异现象}：在评估 Agent 持续改进能力的 \href{https://arxiv.org/abs/2406.08747}{StreamBench} 实验中，研究者发现：\textbf{如果把模型过去"答错"的经验放进 Context 试图让它避免重蹈覆辙，模型的表现反而会显著下降！}就像对人说"不要想白熊"，他反而满脑子都是白熊。目前而言，只给模型提供成功的经验效果最好。

\subsection{2. 压缩 (Compression)}

当互动记录太长，又不想让模型失忆时，可以使用一个专门的模型对历史记录进行摘要。

\keyslide{../key_slides/p4_slide_0195.jpg}{压缩与长期记忆}

\begin{itemize}[leftmargin=2em]
  \item \textbf{递回式压缩}：每互动 100 轮，就将这 100 轮的内容与上一次的摘要进行合并，生成新的摘要（\href{https://arxiv.org/abs/2308.15022}{Recursively Summarizing Enables Long-Term Dialogue Memory}）。太久远的细节会慢慢被遗忘，只留下大纲。
  \item \textbf{对于 Computer Use 尤为重要}：模型与网页互动的过程中充满了"鼠标移至此处、关闭弹窗广告、点击下拉菜单"等垃圾信息。这几百步的操作完全可以被压缩成一句话："已成功预定 A 餐厅，9 月 19 日下午 6 点，10 人"。
\end{itemize}

\keyslide{../key_slides/p4_slide_0200.jpg}{储存到硬碟}

\textbf{防范摘要丢失细节的技巧}：如果不放心摘要会丢掉关键细节，可以把原始细节保存到硬盘（txt 文本）中。在摘要里留下一句话："想要回忆那年夏天的详细故事，请打开 \texttt{xxx.txt} 档案"。这就为模型提供了一条随时读取细节的索引路径。

\subsection{3. 多智能体协作 (Multi-Agent)}

将大任务拆分给多个 Agent 共同完成。

\keyslide{../key_slides/p4_slide_0205.jpg}{Multi-Agent 概念}

从 Context Engineering 的角度来看，Multi-Agent 其实是一种\textbf{卓越的上下文隔离手段}：
\begin{itemize}[leftmargin=2em]
  \item 假设由一个"主 Agent（总召）"负责制定旅行计划。
  \item 它不需要亲自去订餐厅和酒店，而是把"订餐厅"的任务派发给 Agent 1，把"订酒店"的任务派发给 Agent 2。
  \item Agent 1 在订餐厅过程中产生的所有繁琐的网页交互 Context，全都在它自己的窗口里，\textbf{不会污染}主 Agent 的 Context。
  \item 最终主 Agent 得到的 Context 非常清爽："餐厅已定好"、"酒店已定好"。
\end{itemize}

同样，当需要撰写一篇涉及 100 篇论文的 Review 文章时，让一个 Agent 读 100 篇论文肯定会塞爆。但如果派 100 个 Agent 各自读 1 篇并写出摘要，再交给主 Agent 统整，就能完美避开上下文瓶颈。

\keyslide{../key_slides/p4_slide_0210.jpg}{Multi-Agent Performance}

根据 LangChain 等框架的测试，在简单任务上，Single Agent 可能更快更好；但面对高度复杂的任务，Multi-Agent 架构展现出了压倒性的优势。

% ============================================================
\section{总结}
% ============================================================

\begin{itemize}[leftmargin=2em]
  \item \textbf{Context Engineering} 不是什么神奇的"新玄学"，它是管理 AI Agent 大脑缓存的一门务实科学。
  \item 不要盲目迷信超大上下文窗口模型。即便它们"装得下"，它们也"读不懂"、"记不清"。
  \item 熟练运用\textbf{挑选相关资讯（RAG）、压缩无关废话（摘要）、利用团队分工（Multi-Agent隔离）}，才能让你的 AI Agent 稳定、持久、聪明地运行。
\end{itemize}

\end{document}
